% Part: sets-functions-relations
% Chapter: relations
% Section: reflections
%
\documentclass[../../../include/open-logic-section]{subfiles}

\begin{document}

\olfileid{sfr}{rel}{ref}
\olsection{哲学反思}

在 \olref[set]{sec} 中，我们将关系定义为特定的集合。我们应当稍作停顿，简单问一个哲学问题：这样的定义究竟是在做什么？我们极不应当主张我们\emph{发现}了某种形而上学的同一性事实；例如，自然数上的序关系\emph{竟然是}我们在 \olref[set]{sec} 中定义的集合 $R=  \Setabs{\tuple{n,m}}{n, m \in \Nat\text{ and }
n < m}$。理由有三。

首先：在 \olref[set][pai]{wienerkuratowski} 中，我们定义了 $\tuple{a, b} =
\{\{a\}, \{a, b\}\}$。现在考虑定义 $\lVert a,
b\rVert = \{\{b\}, \{a, b\}\} = \tuple{b,a}$。当 $a \neq b$ 时，我们有 $\tuple{a, b} \neq \lVert a,b\rVert$。但我们同样可以把 $\lVert a,b\rVert$ 而非 $\tuple{a,b}$ 视为对有序对的定义。两种定义同样有效。所以现在我们有两个同样有资格“作为”自然数上序关系的候选者，即：
\begin{align*}
		R &= \Setabs{\tuple{n,m}}{n, m \in \Nat \text{ 与 }n < m}\\
		S &= \Setabs{\lVert n,m\rVert}{n, m \in \Nat \text{ 与 }n < m}.
\end{align*}
既然 $R \neq S$，由外延性原理可知，它们不可能\emph{都}与~$\Nat$ 上的序关系同一。但若声称事实上是 $R$ 而非 $S$（或者反过来）\emph{就是}那个序关系，那将纯属任意，因而也有些令人难堪。（这是 Benacerraf（贝纳塞拉夫）在 \citeyear{Benacerraf1965} 中使之著名的一个反对集合论还原主义论证的极简例证。我们稍后会多次回到这个问题。）

其次：如果我们认为\emph{每一个}关系都应被等同于一个集合，那么集合属于关系本身，即 $\in$，也应该是一个特定的集合。事实上，它必须是集合 $\Setabs{\tuple{x,y}}{x \in y}$。但这个集合存在吗？考虑到罗素悖论，声称这样一个集合存在绝非一个平凡的主张。事实上，\oliflabeldef{cumul:::part}{我们在 \olref[cumul][][]{part} 中将发展的集合论将\emph{否定}这个集合的存在。\footnote{提前说明一下原因。用归谬法，假设 $I =
\Setabs{\tuple{x,y}}{x \in y}$ 存在。那么 $\bigcup \bigcup I$ 就是万有集，这与 \olref[sfr][z][sep]{thm:NoUniversalSet} 矛盾。}}{可以以公理化理论的方式严谨地发展集合论，而该理论确实会否定这个集合的存在。}因此，即使某些关系可以被当作集合来处理，集合属于关系也必须是一个特殊情况。

第三：当我们把关系“同一于”集合时，我们曾说可以用 $Rxy$ 来写 $\tuple{x,y} \in R$。只要属于关系“$\in$”被当作一个\emph{谓词}来处理，这就没问题。但如果我们认为“$\in$”代表某个特定的集合，那么表达式“$\tuple{x,y} \in R$”就仅仅由三个代表集合的单称词项组成：“$\tuple{x,y}$”、“$\in$”和“$R$”。而这样一个名字列表，并不比无意义的字符串“杯子 笔筒 桌子”更能表达一个命题。再次，即使某些关系可以作为集合处理，集合属于关系也必须是特殊情况。（这融合了Frege（弗雷格）的“概念\emph{马}”悖论的一个简化版本，以及维特根斯坦（Wittgenstein）曾经对罗素提出的一个著名反驳。）

那么，这对我们有何启示？其实，我们说数上的关系是集合，这本身没有任何\emph{错误}。我们只需理解这话的精神实质。我们并非在陈述一个形而上学的同一性事实。我们只是指出，在某些语境下，我们可以（并且将会）\emph{把}（某些）关系\emph{当作}某些集合来处理。

\end{document}